\section{法、学与寺院：金如何安排北方的共同生活}

\subsection{法令抵达地方之前}

一条法令从中都发出，并不会以同一形状到达所有地方。它要经过抄写、传递、解释和执行；沿途会遇到官员更替、道路阻塞、地方惯例和当事人的抵制。金代法制材料所能首先证明的，是朝廷试图把哪些行为分为可许、可罚、可征或可免；它不能自动证明县衙、军镇和家户已按此处理。正因为如此，法令的反复颁行不是无关紧要的文书堆积，而是制度需要不断同现实摩擦的痕迹。

诉讼的材料保存更为稀少。后出的故事、笔记和地方志有时会保留纠纷的轮廓，却常把人物塑造成道德范型；碑志记下官员断案，通常为了称扬其德。没有一份完整的判牍库，便不能把金代司法写成一套无缝运行的体系。我们仍可看见司法为何重要：土地、债务、婚姻、盗窃、军役和寺院财产都可能需要一个被承认的裁决场所。谁能进入这个场所、谁能负担文书和往返、谁只能在宗族或邻里中解决，决定了国家法律在社会中的实际边界。

\subsection{官员与地方知识}

金廷需要能够读写、核算和处理案牍的官员，因而学校、考试和文书训练并非文化的边缘。它们提供了把中枢命令转为地方表格、账簿和判语的人。可是，官员不只是制度的管道。他们来自家族和地域，有个人的仕途、同僚和地方关系；考课、升迁和战争压力会影响他们究竟重视催税、安民还是自保。把官僚写成全然有效或全然腐败，都不能解释政策为何在某些地点改变、在另一些地点停滞。

金朝采用并调整前代的文官制度，也保留自身军政组织。两者的并置不应被理解为一方逐渐吞没另一方，而是统治者不断处理不同资源与不同权威的办法。军政首领掌握的人员和马匹，州县官掌握的簿册和诉讼，都可能成为中枢不可缺少、又必须防范的力量。一个政权能够把这些关系维持多久，取决于政治信任、财政供给和地方环境，而不是制度名称是否听来整齐。

\subsection{文字并存不是身份统计}

金代的汉文、女真文以及承接辽代传统的多种文字遗存，显示北方并不存在唯一的书写世界。官署、碑刻、宗教和家族纪念选择不同文字，是因为读者、权力和传统不同；同一人或同一地区也可能进入不止一种文字网络。将某种文字直接等同于单一族群、单一政治立场，会把书写实践变成后来的分类工具。

女真文字材料的释读与断代尤须谨慎。残片、拓本、旧藏来源和转写方法都会影响今天的理解。它们有时能够校正汉文史书的官名或人名，有时却只留下难以解释的片段。不能因其稀少而赋予每件遗物过高的代表性，也不能因难读而把金代社会完全交给汉文正史。最好的读法是在具体问题上让不同文字互相照明：某次营建、某位官员、某项宗教捐施或某段道路的称谓，是否在多种材料中留下可比较的痕迹。

\subsection{寺院的门与市场的门}

寺院并非与市场相对的纯净空间。香火、施入、修建、经卷、墓葬和教育都要动用财物和人力；寺院可能成为旅人歇脚、家族纪念和地方交涉的地点。国家可以登记僧道、管理寺产或在危机中调动资源，地方也可能借寺院网络保存声望和互助。这样的关系不能被概括为国家“控制”或宗教“独立”，因为不同寺院、不同阶段和不同地区的力量并不相同。

造像题记和功德碑常以规范的语言赞颂施主、僧人和工程。它们没有提供一张中性社会调查表，却能告诉我们一次营建怎样把姓名、职业、家族与地点固定在石上。题记中缺少的人，也提醒我们物质与书写资源的不平等。一个无法立碑的家庭仍可能在寺院、市场和亲属之间组织生活，只是它的痕迹较难穿过战争、改朝和后来的收藏。

\subsection{文化的成本}

印刷、抄写、教育和礼仪都需要纸张、木料、工匠、时间和相对稳定的保存处。中都或地方城市的文化活动之所以可能，不是因为战争消失，而是因为仍有若干道路、赞助者和机构能够集中资源。反过来，战乱对文化的损害也不只是书籍被焚的戏剧场面：教师和学生离散，纸张和劳力转入军需，寺院和书院的土地收入中断，抄本与档案随迁徙而散失。金末的资料空缺，部分就是这些过程的结果。

因此，评价金代文化不宜只列出诗人、科举和朝廷礼制。应问知识在何处被生产，谁能从中取得资格，哪些机构把文字保存到后世，哪些经验因没有载体而消失。这样的提问也使文化史回到政权转换的中心：一个征服王朝若要在广阔北方治理，必须处理的不只是军队和税粮，还有命令如何被读懂、记忆如何被固定、不同传统如何在同一城市相处。

\subsection{给读者的材料界限}

《金史》的志与列传保存许多制度和人物线索，但由元代官修，且依赖大量已佚的金朝材料；它不应被当作无缝的档案总库。宋方的《建炎以来系年要录》《三朝北盟会编》对于边界、使节和战争不可替代，却从对手政治的角度写作。碑志、题记、都城和寺院考古能使某些地方更清楚，也带有保存和发掘的限制。本章据此辨认问题，而不把它们拼成超过材料承载力的精确数字或动机叙述。
